\documentclass[dvipsnames]{beamer}
\usepackage{comment}
\usepackage{hyperref}

\title{MATH 1190 \& 1210 Meeting for Week 15}
\author{Josh}
\date{20 April 2026}

\newcommand{\transition}[1]{\frame{{}\begin{center}\fbox{#1}\end{center}}}

\begin{document}

\maketitle


\section{Logistics: Pacing}

\transition{Logistics}

\frame{{Logistics: Pacing}

{\bf Where are folks in the lessons right now?}

}

\frame[t]{{Logistics: Upcoming Schedule}

    \begin{itemize}
     \item week 15 of class (20 April - 24 April)
        \begin{itemize}
        \item Application Day 4: Markov Chains\\
        \item {\color{Red} No Preclass!}
        \item Review for final
        \end{itemize}
         \item week 14 of class (27 April - 1 May)
        \begin{itemize}
        \item Review for final
        \end{itemize}
    \end{itemize}

}

\frame[t]{{Logistics: Upcoming Deadlines}

{\bf Assessment-related Logistics}
\begin{itemize}
\item KC4 23 April - CHEMISTRY BLDG ROOM 402 and WARNER HALL ROOM 209
item KC4 23 April
    \begin{itemize}
    \item Proctors: Vadim, Susan, Holden, Michael\\
    \item reassessment targets: {\color{RedOrange}A2},{\color{RedOrange}A3}, {\color{Fuchsia}I1}, {\color{Fuchsia}I2}, 
    \item new targets:{\color{Fuchsia}I3},{\color{Fuchsia}I4},{\color{ProcessBlue}F3}, {\color{Orange}P1}, {\color{Orange}P2} (5 targets!)
    \item Thursday, 23 April, 7:00 PM; 90 minutes
    \end{itemize}
\item Grading for KC4 needs to be finished by Wednesday May 29
\item KC4 regrade requests will open that night and close the next night - be sure to tell your students!
\item KC4 Reflection? Do they have time?

\end{itemize}

}

\frame[t]{{Logistics: Upcoming Deadlines}

{\bf Assessment-related Logistics}
\begin{itemize}

item KC4 23 April
    \begin{itemize}
    \item new targets:
    \begin{itemize}
    \item {\color{Fuchsia}I3}: Michael and Josh\\
    \item {\color{Fuchsia}I4}: Everyone!\\
    \item {\color{ProcessBlue}F3}: Jim B. and Vadim\\
    \item {\color{Orange}P1}: Susan and Jim R.\\
    \item {\color{Orange}P2}: Aaratrick and Holden 
    \end{itemize}
    \end{itemize}
\item Final Exam - 1 May 7:00pm - 10:00pm 
\begin{itemize}
    \item Ridley GOO8: 
    \item Physics 238: 
    \item ``Physics 338": 
\end{itemize}
\end{itemize}

}


\section{Content}

\transition{Content}


%------------------------------------------------


%%%%%%%%%%%%%%%%%%%%%%%%%%%%%%%%%%%%%%
\begin{frame}{Activity Overview}
\begin{itemize}
    \item Students will explore a real-world system using a Markov Chain.
    \begin{itemize}
        \item Interpret a transition diagram
        \item Build a transition matrix
        \item Track how a system evolves over time
        \item Make predictions about long-term behavior
    \end{itemize}
    \item Work with a partner and justify reasoning.
\end{itemize}
\end{frame}

%%%%%%%%%%%%%%%%%%%%%%%%%%%%%%%%%%%%%%
\begin{frame}{Step 1: Understand the Scenario}
\begin{itemize}
    \item Identify:
    \begin{itemize}
        \item The \textbf{states}
        \item What it means to move between states
    \end{itemize}
    \begin{itemize}
        \item What do the probabilities represent?
    \end{itemize}
\end{itemize}
\end{frame}

%%%%%%%%%%%%%%%%%%%%%%%%%%%%%%%%%%%%%%
\begin{frame}{Step 2: Build the Transition Matrix}
\begin{itemize}
    \item Students will use a diagram to construct matrix $P$.
    \begin{itemize}
        \item Each column must sum to 1
        \item Entries represent transition probabilities
    \end{itemize}
    \item Idea: Have them check their matrix with another group.
\end{itemize}
\end{frame}

%%%%%%%%%%%%%%%%%%%%%%%%%%%%%%%%%%%%%%
\begin{frame}{Step 3: Predict Short-Term Behavior}
\begin{itemize}
    \item Start with an initial state vector $S$.
    \item Compute:
    \[
    S_1 = P S, \quad S_2 = P^2 S
    \]
    \item Interpret your results:
    \begin{itemize}
        \item What do these vectors represent?
    \end{itemize}
\end{itemize}
\end{frame}

%%%%%%%%%%%%%%%%%%%%%%%%%%%%%%%%%%%%%%
\begin{frame}{Step 4: Long-Term Behavior}
\begin{itemize}
    \item Continue applying the matrix (or reason conceptually).
    \item Answer:
    \begin{itemize}
        \item Does the system stabilize?
        \item What do the long-term values mean?
    \end{itemize}
    \item Explain your reasoning.
\end{itemize}
\end{frame}

%------------------------------------------------

%%%%%%%%%%%%%%%%%%%%%%%%%%%%%%%%%%%%%%%%%%%%%%%%%%%%%%%%%%%%%%%%%%%%%%%%%%%%%%%%%%%%%%%%%%%%%%%%%%%%%%
%%%%%%%%%%%%%%%%%%%%%%%%%%%%%%%%%%%%%%%%%%%%%%%%%%%%%%%%%%%%%%%%%%%%%%%%%%%%%%%%%%%%%%%%%%%%%%%%%%%%%%

\section{Pedagogy}

\transition{Pedagogy}


\transition{Active Learning}
\transition{And Now - An Active Learning Discussion with Vadim!}
%%%%%%%%%%%%%%%%%%%%%%%%%%%%%%%%%%%%%%%%%%%%%%%%%%%%%%%%%%%%%%%%%%%%%%%%%%%%%%%%%%%%%%%%%%%%%%%%%%%%%%

%------------------------------------------------



%------------------------------------------------

%%%%%%%%%%%%%%%%%%%%%%%%%%%%%%%%%%%%%%
\begin{frame}{Final Exam Review in a Growth-Based Setting}
\begin{itemize}
    \item Traditional final exams are often:
    \begin{itemize}
        \item High-stakes
        \item One-shot assessments
        \item Focused on point accumulation
    \end{itemize}
    
    \item Growth-based grading (Standards-Based Grading / specifications grading):
    \begin{itemize}
        \item Emphasizes mastery over time
        \item Encourages revision and persistence
    \end{itemize}
    
    \item Key Question:
    \begin{center}
        \textit{How should final exam preparation change under this model?}
    \end{center}
\end{itemize}
\end{frame}

%%%%%%%%%%%%%%%%%%%%%%%%%%%%%%%%%%%%%%
\begin{frame}{Relevant Literature}
\small
\begin{itemize}
    \item \textbf{Student anxiety and Standards-Based Grading:}
    \begin{itemize}
        \item \href{https://link.springer.com/article/10.1007/s10755-019-09489-3}{Student Anxiety in Standards-Based Grading in Mathematics Courses}
    \end{itemize}
    
    \item \textbf{Growth mindset and mastery:}
    \begin{itemize}
        \item \href{https://peer.asee.org/fostering-growth-mindsets-implementing-standards-based-grading-in-college-algebra}{Fostering Growth Mindsets: Implementing Standards-Based Grading in College Algebra}
    \end{itemize}
    
    \item \textbf{Specifications grading in mathematics:}
    \begin{itemize}
        \item \href{https://www.tandfonline.com/doi/full/10.1080/0020739X.2025.2519130}{Beyond Traditional Grading: Specifications Grading in Mathematical Biology}
    \end{itemize}
    
    \item \textbf{Practical teaching guidance:}
    \begin{itemize}
        \item \href{https://www.umass.edu/ctl/how-do-i-design-my-grading-system-growth-mind}{UMass CTL: Designing Grading Systems with Growth in Mind}
    \end{itemize}
\end{itemize}
\end{frame}

%%%%%%%%%%%%%%%%%%%%%%%%%%%%%%%%%%%%%%
\begin{frame}{Key Idea 1: Diagnostic, Not Performative}
\begin{itemize}
    \item Shift from:
    \begin{itemize}
        \item ``Can you get this right right now?''
    \end{itemize}
    
    \item To:
    \begin{itemize}
        \item ``What do you still need to learn?''
    \end{itemize}
    
    \item Strategies:
    \begin{itemize}
        \item Important ideas checklists
        \item Targeted practice by sub-concepts
    \end{itemize}
\end{itemize}
\end{frame}

%%%%%%%%%%%%%%%%%%%%%%%%%%%%%%%%%%%%%%
\begin{frame}{Key Idea 2: Build in Revision Cycles}
\begin{itemize}
    \item Learning happens through iteration
    \item Use:
    \begin{itemize}
        \item Error analysis: ``What is wrong with this solution?"
        \item Revising incorrect solutions
        \item Comparing multiple approaches
    \end{itemize}
    
    \item Emphasize:
    \begin{itemize}
        \item Process over final answer
    \end{itemize}
\end{itemize}
\end{frame}

%%%%%%%%%%%%%%%%%%%%%%%%%%%%%%%%%%%%%%
\begin{frame}{Key Idea 3: Student Self-Assessment}
\begin{itemize}
    \item Students identify:
    \begin{itemize}
        \item Strengths
        \item Areas for growth
    \end{itemize}
    
    \item Example prompts:
    \begin{itemize}
        \item ``Which learning targets are you least confident in?''
        \item ``What is your plan to improve?''
    \end{itemize}
    
    \item Builds metacognition and ownership
\end{itemize}
\end{frame}

%%%%%%%%%%%%%%%%%%%%%%%%%%%%%%%%%%%%%%
\begin{frame}{Key Idea 4: Reduce One-Shot Pressure}
\begin{itemize}
    \item Frame the final exam as:
    \begin{itemize}
        \item Another opportunity to demonstrate learning
        \item Opportunity to replace previous learning targets you may have not mastered
    \end{itemize}
    
    \item Helps reduce anxiety and improve performance
\end{itemize}
\end{frame}

%%%%%%%%%%%%%%%%%%%%%%%%%%%%%%%%%%%%%%
\begin{frame}{Key Idea 5: Depth Over Coverage}
\begin{itemize}
    \item Avoid:
    \begin{itemize}
        \item Large lists of disconnected review problems
    \end{itemize}
    
    \item Focus on:
    \begin{itemize}
        \item Explaining reasoning
        \item Making connections
        \item Analyzing common mistakes from previous exams
        \item Look at previous rubrics and their main ideas
    \end{itemize}
    
\end{itemize}
\end{frame}

%%%%%%%%%%%%%%%%%%%%%%%%%%%%%%%%%%%%%%
\begin{frame}{Reflection Questions}
\begin{itemize}
    
    \item What would it look like to make your review process more diagnostic?
    
    \item How could we incorporate revision or error analysis into exam preparation? 
    
    \item In what ways do your students currently engage in self-assessment?

\end{itemize}
\end{frame}

%%%%%%%%%%%%%%%%%%%%%%%%%%%%%%%%%%%%%%
\begin{frame}{Discussion Questions }
\begin{itemize}

\item \textbf{What does “alignment” look like between:}
\begin{itemize}
    \item Standards/learning goals
    \item Review materials
    \item The final exam itself?
\end{itemize}

\vspace{0.4cm}

\item \textbf{Compare two models for review in a coordinated course:}
\begin{itemize}
    \item Centralized review materials (same for all sections)
    \item Instructor-driven, standards-based review
\end{itemize}
\textit{What are the benefits and limitations of each?}

\vspace{0.4cm}


\item \textbf{What constraints matter most in this course?}
\begin{itemize}
    \item Consistency across sections?
    \item Instructor autonomy?
    \item Student experience?
\end{itemize}

\end{itemize}
\end{frame}
%------------------------------------------------

\end{document}


